\section{Introduction}\label{sec:intro}

A repair centre for medical imaging equipment receives field-replaceable units from hospitals: power supplies,
amplifiers, control boards, detector modules. Most are repaired and returned to service. The others are found beyond
economic repair and are \emph{scrapped}. At the site that motivated this work, scrapped units leave whole, through a
bulk waste contract that produces no revenue. Yet a scrapped unit is rarely a worthless one. It failed for one cause,
and the components that did not cause the failure are often sound: they can be extracted, tested and \emph{reused}
in later repairs, where they replace parts that are costly and slow to procure. What cannot be reused can be
\emph{sorted} by material and sold by the kilogram to specialised recyclers, who pay for printed circuit boards,
copper or aluminium many times what a mixed stream is worth. The alternative to the bulk contract is therefore a small
teardown line: open the unit, take out what is worth taking, test it, sort the rest.

Whether to install such a line, and how to run it, looks like a textbook disposition problem. It is not, and the
reason is the subject of this paper. A disposition model needs, for every component of every unit type, the
probability that it is still good, the time to extract and test it, its value if reused, its mass, and the price a
recycler pays for its material class. At the start of the project almost none of these were known.
\begin{itemize}
\item \emph{Yields are unknown.} Nobody has torn down scrapped units, so nobody knows which share of their components
  would pass a test. Consumption records of repairs give a prior, biased in directions that can be named but not
  quantified without teardown data.
\item \emph{Arrivals are irregular.} Scrapped units arrive a few per week per type, with many empty weeks, while
  operator hours are booked in advance.
\item \emph{Compositions are missing.} The bill of materials gives a company part number and a free-text description.
  The material class of a component, which fixes its recycling price and its regulatory route, is recorded nowhere.
  It can be found: by a rule on the description, by analogy with a known part, in a manufacturer's material
  declaration, by reading a data sheet. Each of these has a cost and can be wrong, and there are thousands of parts.
\item \emph{The data are private.} Part numbers and descriptions of in-house designs cannot leave the company, which
  excludes sending them to a hosted language model, the tool that would otherwise read descriptions and data sheets
  at scale.
\end{itemize}

The system we describe treats these difficulties as one problem with three layers. The \emph{disposition} layer
decides, for each unit type, whether to open it and what to do with each component, under a labour budget. The
\emph{learning} layer estimates yields and arrivals while decisions are being taken. The \emph{knowledge} layer
decides which missing facts are worth acquiring, by which means, and when to stop. The layers are coupled by a single
quantity. The dual price of labour computed by the disposition layer prices every unit and every component; the same
optimisation model, evaluated with and without a fact, gives the value of knowing it; and that value is the budget
the knowledge layer may spend to acquire it. The result is a system in which a language model is called on a
component only when the answer can pay for the call, and a human is asked only where the stake exceeds the cost of
asking.

\paragraph{Contributions.} The paper makes four contributions.

\emph{(1) Structure of the disposition problem.} The problem is the linear relaxation of a multiple-choice knapsack
over unit types, with one item per plan (a disposition for each component) and a set-up that depends on how far the
unit is opened. We show that, at a given labour price, the best plan of a unit is found by pricing its components
independently for each access level (Proposition~\ref{prop:pricing}), which removes the enumeration of plans and makes
units of a thousand components as easy as units of five. The optimal labour price coordinates all decisions
(Propositions~\ref{prop:dual} and~\ref{prop:split}), and gives the operator a threshold on the observed good rate of
a component above which it is worth extracting and testing (Proposition~\ref{prop:threshold}). With irregular
arrivals, the gap between the model on the average flow and the model replayed period by period is an upper bound on
what \emph{any} policy that stores units, forecasts arrivals or is trained by reinforcement learning can gain over
the myopic policy (Proposition~\ref{prop:fluid}). The bound is computable from the history, so the question ``should
we build a dynamic policy?'' is answered by a number before anything is built.

\emph{(2) A budgeted knowledge layer.} We define the value of establishing the material class of a component as the
difference between two optimal values of the disposition model, and bound it by a quantity computable from the
bill of materials alone (Proposition~\ref{prop:class}). The acquisition chain spends on a component only within that
bound. We prove that the chain never spends on a component more than its class can be worth
(Lemma~\ref{lem:budget}), that an acquisition excluded by the screen on a component not yet paid for has negative
expected net value (Proposition~\ref{prop:screen}), that acceptance without review is the cheaper decision whenever the value at risk
is below the cost of a review (Lemma~\ref{lem:accept}), that the fusion of sources identifies the true class at an
exponential rate whenever a source's reliability differs from chance (Proposition~\ref{prop:fusion}), and that
reliabilities estimated from human validations are unbiased only if the validations are stratified by the reason the
component was validated (Proposition~\ref{prop:strata}). In this layer a language model is a source with a learned
reliability, never ground truth.

\emph{(3) Agents on a self-hosted open model.} The whole system is driven by three agents (a pilot, a writer, a
reviewer) that run on an open-weight language model hosted on the data owner's machine. They can only call a closed
list of the system's own functions; every number of the decision note they write is checked by a program against
the computed results (Lemma~\ref{lem:grounding}); without a model, a fixed plan produces the same note from a
template. The instruction given to the model for reading part descriptions is itself improved from the human
validations by reflective prompt evolution, with a call budget and a held-out test
(Lemma~\ref{lem:heldout}). No weight is trained and no record leaves the site.

\emph{(4) An open, tested implementation and its application.} All of the above is implemented in a single
package with a test suite that checks, among other things, the optimum against an exact linear-programming solver,
every proposition on random instances, and the confidentiality of outgoing queries.
Section~\ref{sec:experiments} measures the algorithms on synthetic instances, and Section~\ref{sec:case} reports
the application to the repair centre.

\paragraph{Status of the evidence.} We separate three kinds of statement. Mathematical results are proved in
Appendix~\ref{app:proofs} and checked numerically. Measurements of the algorithms use synthetic instances and say
nothing about the site. Statements about the site come from its data. \ifnum\DonneesDemo=1\relax In this draft the
case study runs on the fictional demonstration data shipped with the software, and every sentence that depends on
field results is marked as still to be written.\fi{} We make no claim that is not of one of these three kinds.

\paragraph{Organisation.} Section~\ref{sec:literature} reviews the literature. Section~\ref{sec:model} states the
disposition model and Section~\ref{sec:structure} its structure. Section~\ref{sec:learning} treats unknown yields and
irregular arrivals, Section~\ref{sec:knowledge} the acquisition of missing compositions, and Section~\ref{sec:system}
the agents and the deployment. Sections~\ref{sec:experiments} and~\ref{sec:case} give the computational study and
the case study. Section~\ref{sec:discussion} discusses managerial insights, limits and the transfer to production.
